\bookchapter{Глава 11: Фотография и управление снимками}

\chapterintro{Скажу честно: большинство людей снимают на телефон далеко не на пределе его возможностей. Дело не в ваших навыках, а в том, что вы не замечаете множества мелких настроек и привычек. В этой главе — только то, что можно применить прямо сейчас.}

\booksection{Допускаете ли вы эти ошибки?}

Давайте разберём самые частые проблемы — проверьте, сколько из них ваши:

\begin{itemize}
  \item \textbf{Грязный объектив}: это действительно самая распространённая проблема. Достаёте телефон — палец попадает на объектив, снимок получается мутным, и вы думаете, что телефон сломался. Выработайте привычку: перед съёмкой протирайте объектив краем одежды.
  \item \textbf{Никогда не настраиваете экспозицию}: в обычном режиме касание экрана фокусирует, но если после касания провести пальцем вверх-вниз — можно регулировать яркость. Многие снимки получаются слишком тёмными или слишком светлыми именно потому, что люди не знают об этом жесте.
  \item \textbf{Вспышка всегда включена}: светодиодная вспышка телефона даёт ужасный результат — жёсткий, холодный свет, делающий людей похожими на призраков. Выключайте её всегда, кроме случаев кромешной тьмы. В условиях слабого освещения ночной режим даёт гораздо лучший результат, чем вспышка.
  \item \textbf{Приближение щипком двумя пальцами}: масштабирование щипком на экране — это цифровой зум, по сути — обрезка и растягивание картинки, качество заметно падает. Если нужно снять удалённый объект, используйте кнопки зума в видоискателе (0.5x/1x/2x/3x). 0.5x (сверхширик) и 1x (основная камера) — это нативные оптические объективы без потери качества; 2x/3x — чаще всего «зум оптического качества», достигаемый кропом с высокого разрешения сенсора, качество близко к оптическому, но это не отдельный объектив. Зум выше 3x — обычно гибридный или цифровой, с заметным падением качества.
\end{itemize}

\booksection{Несколько приёмов для мгновенного повышения выхода удачных кадров}

Никакой теории — попробуйте раз и увидите результат:

\begin{itemize}
  \item \textbf{Правило третей + сетка}: включите «Сетку» (направляющие) в настройках камеры. При съёмке помещайте главный объект на пересечение линий — и кадр сразу станет «правильным». Это самое простое и эффективное правило в фотографии, никакой магии — просто соответствует визуальным привычкам человеческого глаза.
  \item \textbf{Ищите «направляющие линии»}: снимая пейзаж или городскую сцену, обратите внимание на естественные линии в кадре — дорогу, перила, стену, ряд деревьев — и пусть они ведут от края кадра к главному объекту. Освоив этот приём, вы будете получать снимки, которые прохожие примут за работу человека, учившегося фотографии.
  \item \textbf{Не снимайте всегда стоя}: большинство людей снимает с высоты собственных глаз — отсюда однообразие ракурса. Попробуйте присесть и снять с низкой точки или найти возвышение и снять сверху — одно и то же место будет выглядеть совершенно иначе. Снимая детей и животных, особенно важно присесть — съёмка на уровне глаз даёт естественный результат.
  \item \textbf{Свет — это главное}: портрет под полуденным солнцем даёт жёсткие, уродливые тени на лице. Лучший свет — «золотые часы» на рассвете и закате: косые лучи, тёплый оттенок — всё выглядит прекрасно. Пасмурный свет пусть и не эффектен, зато равномерен и мягок — портреты в такую погоду получаются даже лучше, чем на ярком солнце.
  \item \textbf{Как пользоваться ночным режимом}: включив ночной режим, телефон делает серию снимков и объединяет их. В этот момент критически важно — руки не должны дрожать. Прислонитесь к стене, положите телефон на стол, в крайнем случае задержите дыхание. Ночной режим требует неподвижности в течение 2-4 секунд — малейшее дрожание смажет кадр.
\end{itemize}

\booksection{Портретный режим и профессиональный режим}

\begin{itemize}
  \item \textbf{Портретный режим}: создаёт эффект размытого фона при чётком объекте (малая глубина резкости). Из-за физически маленького объектива телефона реальное размытие ограничено, поэтому эффект достигается алгоритмами AI — проще говоря, AI «вырезает» человека и размывает фон. Лучший результат — на расстоянии 1-2 метра; слишком далеко или слишком близко — алгоритм чаще ошибается. И ещё: при портретной съёмке избегайте слишком сложных линий на фоне, иначе AI «провалится» на границах волос.
  \item \textbf{Профессиональный режим}: если хотите полностью контролировать картинку, попробуйте профессиональный режим. Ключевые параметры: чем ниже ISO — тем чище картинка, но темнее; чем длиннее выдержка — тем больше света, но дрожание рук смажет кадр; баланс белого управляет теплотой/холодностью цвета — при неправильной настройке фото уйдёт в синеву или желтизну. Обычному человеку не нужно постоянно снимать в профессиональном режиме, но понимание этих параметров хотя бы даст вам представление о том, что делает камера.
\end{itemize}

\booksection{Постобработка: не переусердствуйте, достаточно лишь нужного}

Многие неправильно понимают ретушь, считая её «ненастоящей». На самом деле снимок с камеры и так отличается от того, что видит глаз; умеренная коррекция лишь восстанавливает увиденную вами картину, а не подделывает её.

\begin{itemize}
  \item \textbf{Snapseed}: бесплатный инструмент от Google, профессиональный и без рекламы. Мой повседневный процесс: «Коррекция изображения» — подтянуть яркость и контраст «Усиление деталей» — немного резкости «Баланс белого» — исправить цветовой сдвиг «Кадрирование» — довести композицию. Функция «Локально» в Snapseed очень удобна — позволяет править только выбранную область, не затрагивая всё изображение.
  \item \textbf{Xingtu / Meitu Xiuxiu}: эти два хороши для портретной ретуши. Но, на мой взгляд, есть принцип: фильтры красоты — не более 30\%, иначе выглядит фальшиво, кожа как пластик. Естественность красивее совершенства.
  \item \textbf{Lightroom для телефона}: если вы требовательны к цвету — это самый профессиональный инструмент. Кривые, HSL-коррекция — всё есть. Но для большинства людей это как из пушки по воробьям.
  \item \textbf{Рекомендация по рабочему процессу}: я привык сначала снимать хороший исходник, а обрабатывать уже дома на большом экране компьютера или iPad — экран телефона слишком мал, да и цветопередача не идеальна. Если приходится править на телефоне — выкрутите яркость экрана выше среднего и не редактируйте в режиме энергосбережения.
\end{itemize}

\booksection{Слишком много фото? Стратегия организации}

Тысячи снимков в галерее, и чтобы найти конкретный, приходится листать бесконечно — знакомая ситуация. Несколько практичных методов:

\begin{itemize}
  \item \textbf{Используйте поиск}: сейчас и iPhone, и большинство Android-телефонов поддерживают интеллектуальный поиск по галерее. Можно искать по ключевым словам: «кот», «море», «паспорт», «цветы» — система AI распознает содержимое снимков и найдёт нужное. Можно искать по дате и месту, например, «май прошлого года», «Пекин». Об этой функции знают немногие, а она действительно полезна.
  \item \textbf{Регулярно удаляйте мусор}: скриншоты, дубликаты, смазанные кадры — вот главные пожиратели памяти. В iOS Галерея автоматически объединяет «Дубликаты»; на Android можно зайти в файловый менеджер и найти папку «screenshots» для массового удаления ненужных скриншотов. Я трачу пять минут в месяц на очистку — и освобождаю несколько гигабайт.
  \item \textbf{Создавайте альбомы, не полагайтесь только на сортировку по дате}: сортировка по времени по умолчанию — это нормально, но для таких межвременных тем, как «Рост ребёнка», «Фото из путешествий», «Важные документы», стоит вручную создать альбомы. «Общие альбомы» Apple позволяют всей семье добавлять туда фото — один создал, все видят.
  \item \textbf{Резервное копирование — это минимум}: фотографии — самые незаменимые данные. Файлы можно восстановить, переписка — не критична, но детские фото и видео с пожилыми родственниками, потерянные навсегда, — это действительно потеря. Обязательно включите облачное резервное копирование. iCloud Фото — самый беспроблемный вариант; на Android у каждого бренда свой облачный сервис, либо используйте Baidu Wangpan / Alibaba Cloud Drive для автозагрузки. Принцип: минимум одна облачная копия, не кладите все яйца в одну корзину.
\end{itemize}

\booksection{Скриншоты и запись экрана}

\begin{itemize}
  \item \textbf{Скриншот}: одновременно нажмите кнопку питания + кнопку увеличения громкости — универсальный метод для всех телефонов. Многие Android-телефоны также поддерживают смахивание тремя пальцами вниз или двойное постукивание костяшкой пальца. Длинный скриншот (прокручиваемый) хорошо реализован на Android — после снятия нажмите кнопку «Прокрутка», и телефон автоматически пролистнёт содержимое и склеит в одно длинное изображение. Подходит для истории чатов и веб-страниц.
  \item \textbf{Запись экрана}: добавьте «Запись экрана» в Пункт управления, нажмите для старта, ещё раз — для остановки. Перед записью долгое нажатие на кнопку записи позволяет выбрать, записывать ли звук с микрофона — включите, если делаете обучающее видео с голосом, выключите, если только записываете действия на экране. Готовое видео автоматически сохраняется в галерею.
\end{itemize}

\begin{tcbpractice}{Потренируйтесь: полный цикл — съёмка + обработка + организация}
\begin{enumerate}
  \item Включите «Сетку» в настройках камеры, найдите место с естественным светом и сделайте три снимка (не забудьте протереть объектив и настроить экспозицию).
  \item Выберите лучший кадр и сделайте базовую коррекцию в Snapseed (яркость, контраст, кадрирование).
  \item Создайте для этого снимка отдельный альбом с названием «Мои работы».
  \item Проверьте, включено ли облачное резервное копирование — убедитесь, что фото уже сохранено в облаке.
\end{enumerate}

\end{tcbpractice}
